\cotefile{57}
\scanpage{193}
\uline{Remarques}. 1) On notera que la démonstration
montre que si $\dim X=1$, alors $\mathrm{Pic}(\widehat E)\unsure{-|\to}\mathrm{Pic}(\deleted{\unsure{E|\widehat E}}X)$
a un noyau qui est connexe\inserted{ et $\unsure{=|\simeq}$ unipotent}, \deleted{\unsure{le|la}} donc
$\mathrm{Pic}(S-\unsure{\bullet|o|a})\unsure{-|\to}\mathrm{Pic}(\widehat S-\bullet)$ induit un isom sur les
groupes \uline{mod composantes connexes} de l'identité
[qui dans ce cas est fini, isom. à $\mathbb Z/d\mathbb Z$ où
$d$ est le degré projectif de $X$]\inserted{ et $=$ modulo partie unipotente du \unsure{dernier.|2\textsuperscript{d} terme.}} Le $=$ résultat
semble valable si $X$ est de dim quelconque,
lorsque \deleted{\unsure{on se met|on se}} la caractéristique \deleted{résiduelle}
$p$ de $\deleted{\unsure{k}}k$ est $>0$, et qu'on raisonne
\uline{mod} \uline{$p$-groupes}\inserted{ \uline{de $p$-torsion}}; i.e. le groupe de N.S. local
semble le $=$ pour $S$ et $\widehat S$, modulo des
$p$-groupes [ext. d'un $p$-groupe fini par un
groupe alg. unipotent sur $k$]; \unsure{utiliser|utilise|utilisant} le
fait que les obstructions dans des
% slanted
\margin{Regarder aussi le cas de la car. résiduelle nulle \ldots\,.}
$H^2(X,\underline O_X(n))$ qu'on rencontre sont nulles
pour $n$ grand, et sont annulés par $p$ \ldots\,. 
Ceci est en accord avec la théorie des
groupes \unsure{fondamentx|fondamentaux} \ldots\,. 
% archive sheet number: 192
